\section{Life without Logos}

\begin{quotex}
We may regard the present state of the universe as the effect of its past and the cause of its future. An intellect which at a certain moment would know all forces that set nature in motion, and all positions of all items of which nature is composed, if this intellect were also vast enough to submit these data to analysis, it would embrace in a single formula the movements of the greatest bodies of the universe and those of the tiniest atom; for such an intellect nothing would be uncertain and the future just like the past would be present before its eyes.
\flright{\textsc{Pierre Simon Laplace}, \emph{A Philosophical Essay on Probabilities}}
\end{quotex}


The dream of modernism and its handmaid science is to explain the \textbf{world}, \textbf{life}, and \textbf{thought} without recourse to anything transcendent. In other words, intelligibility, or logos, is claimed to be immanent in space and time. The \textbf{Principle of Sufficient Reason} states that everything has its own reason for being; thus rationality, or reasonableness, means to know the reason or cause for whatever exists. As was shown in Intelligibility, Ideology and Occult Scientific Theories\footnote{\url{http://www.gornahoor.net/?p=1104}}, the ultimate explanations of science involve uncertainty, indeterminateness, randomness, or unconsciousness. Thus, Science is inherently irrational so its promise of a rational theory of everything has, despite its many local successes, been unfulfilled.

\paragraph{World}


Our knowledge of the world is limited, in principle, by the Heisenberg Uncertainty Principle\footnote{\url{https://en.wikipedia.org/wiki/Heisenberg\_uncertainty\_principle}}. This proves that it is not just difficult, but actually impossible, to know simultaneously both the position and the motion of any particle. Hence the Laplacean project is stillborn.

\paragraph{Life}


The fundamental principle of Darwinian evolution is random variation, which is fundamentally the lack of a rational explanation that would explain in any detail the transition from one species to another. Nevertheless, proponents of evolution subtly, or not so subtly, sneak in teleological explanations. One example is the search for the so-called missing link\footnote{\url{https://en.wikipedia.org/wiki/Missing\_link}}. This implicitly assumes a direction to evolution, from the lower to the higher. However, if evolution is truly random, it seems curious that evolution never runs backwards; for example, a human group could devolve back to apes, or a frog into a fish.

It also tacitly assumes the existence of an essence that changes its form over time. Thus the missing link is said to evolve from one form to the next. In actuality, the missing link is eliminated as unfit for survival, in a sense a ``mistake'' or evolutionary dead end, and replaced by another species.

\paragraph{Thought}


In accounting for the existence of thought, psychology is led to unconscious psychic forces, or even bioelectrochemical processes, that determine the thoughts, and hence the actions, of man. No one can control his thoughts, nor even predict what he will think of five minutes from now. Rather than a ``rational'' animal and master of himself, man is inherently irrational and subject to forces of which he is ultimately unaware.

\paragraph{Summary}


Notice what this argument shows.

\begin{itemize}


\item First of all, it is not anti-science, nor does it dispute the legitimate discoveries and theories of science.

\item Second, the point is not that the scientific enterprise is not as objective and impartial as claimed. This is obvious, not just to a Feyerabend or a Kuhn, but to any observer.

\item Third, it is not a question of the incomplete state of science in which new explanations could, in principle, be discovered in the future.
\end{itemize}


No, the point is that science, insofar as it claims to be profane and independent of transcendence, is fundamentally and irremediably irrational in the sense that it is ignorant of the sufficient reason for the appearance of life, thought, and the world itself. \emph{This is not simply due to the lack of knowledge, but rather to the lack of the very possibility of knowledge.}


\flrightit{Posted on 2010-12-28 by Cologero}

